% Generated draft: requires mathematical and source review.
\documentclass[11pt]{article}
\usepackage[margin=25mm]{geometry}
\usepackage{fontspec}
\setmainfont{Latin Modern Roman}
\usepackage{amsmath,amssymb,mathtools,unicode-math}
\setmathfont{Latin Modern Math}
\usepackage{enumitem,needspace,xurl,hyperref}
\setlength{\parindent}{0pt}
\setlength{\parskip}{5pt}
\setlength{\emergencystretch}{3em}
\newcommand{\N}{\mathbb{N}}
\newcommand{\Z}{\mathbb{Z}}
\newcommand{\Q}{\mathbb{Q}}
\newcommand{\R}{\mathbb{R}}
\newcommand{\C}{\mathbb{C}}
\newcommand{\F}{\mathbb{F}}
\newcommand{\A}{\mathbb{A}}
\newcommand{\PP}{\mathbb{P}}
\newcommand{\E}{\mathbb{E}}
\newcommand{\eps}{\varepsilon}
\newcommand{\dd}{\,\mathrm{d}}
\begin{document}
\section*{UnsolvedMath ID: 20000815}
\textbf{Source problem code:} catalogue entry
\subsection{Definitions and mathematical statement}
The source asks for a geometric/algebraic description of generic points of irreducible components of \(\operatorname{Hilb}^d(\A^n)\), without specifying the ground field or desired classification. Use an algebraically closed field \(k\) of characteristic zero for the supplied explicit constructions. A point corresponds to a colength-\(d\) ideal \(I\subset k[x_1,\ldots,x_n]\), or the finite algebra \(A=k[x_1,\ldots,x_n]/I\). Its decomposition into local Artin factors \(A=\prod_i A_i\) is the disjoint-support decomposition of the associated finite scheme. Call a component elementary when its general member has exactly one support point; the support point may move in \(\A^n\).

The aid's structural theorem says every irreducible component \(C\) has a unique multiset of elementary component factors \(E_1,\ldots,E_s\), with lengths \(d_i\) summing to \(d\). On a dense open subset a member is the disjoint union of general members of these factors, with distinct supports. Repeated factors and permutations are interpreted through this multiset. The disjoint-union map is considered on the open product where supports are disjoint; its local deformation structure gives the reported factorwise invariants.

For a component \(C\), let \(t(C)\) be the generic tangent dimension \(\dim_k\operatorname{Hom}(I,A)\), evaluated at a sufficiently general closed point. Define its dimension excess \(e(C)=\dim C-nd\) and tangent defect \(\delta(C)=t(C)-\dim C\). The asserted formulas are
\[
 \dim C=\sum_i\dim E_i,\qquad t(C)=\sum_i t(E_i),\qquad
 e(C)=\sum_i(\dim E_i-nd_i),\qquad
 \delta(C)=\sum_i\delta(E_i).
\]
Generic reducedness is equivalent to that of all elementary factors in this characteristic-zero setting. A general member is smoothable, meaning a flat limit of \(d\) distinct reduced points, exactly when each of its local factors is smoothable. These statements reduce the generic-point problem to elementary components; they do not classify all elementary factors.

For the aid's explicit concatenations, use known elementary blocks with local Hilbert functions \((1,4,3)\), \((1,5,3)\), and \((1,6,3)\), respectively of lengths \(8,9,10\), in ambient dimensions at least \(4,5,6\). The local Hilbert function is \(h_i=\dim_k\mathfrak m^i/\mathfrak m^{i+1}\). The reported component dimensions for these embedded blocks are \(8n-7\), \(9n-4\), and \(10n\). Taking \(a,b,c\) blocks of these three types and \(s\) single reduced points gives a component \(C_{a,b,c,s}\) with
\[
 a,b,c,s\geq0,\qquad d=8a+9b+10c+s,\qquad
 \dim C_{a,b,c,s}=nd-7a-4b.
\]
For \(n=4\) require \(b=c=0\); for \(n=5\) require \(c=0\); for \(n\geq6\) all are allowed. The support-length and local-Hilbert-function multiset distinguishes these constructed components. If \(c_n(d)\) is the number of irreducible components of \(\operatorname{Hilb}^d(\A_k^n)\), the resulting lower bounds are
\[
 c_4(d)\geq[t^d]\frac1{(1-t)(1-t^8)},\qquad
 c_5(d)\geq[t^d]\frac1{(1-t)(1-t^8)(1-t^9)},
\]
\[
 c_n(d)\geq[t^d]\frac1{(1-t)(1-t^8)(1-t^9)(1-t^{10})}\qquad(n\geq6).
\]
These are lower bounds from selected blocks, not complete counts or a generic-ideal description of every component. The underlying block ideals, tangent calculations, and references are not supplied in the record.
\subsection{Short English statement}
The supplied framework describes a general finite scheme as a disjoint union of elementary one-support pieces, with component and tangent invariants adding across the pieces. Three reported elementary blocks give explicit component families and counting lower bounds. The remaining classification of elementary components is not solved by this decomposition.
\subsection{Sources}
\begin{itemize}
\item Local registry: permanent UnsolvedMath ID \(20000815\), exact public code \texttt{catalogue entry}; \texttt{aim-arithmetic-geometry-notes.json}, \texttt{notes[60]}.
\item AIM workshop \emph{Components of Hilbert schemes}, Problem 14: \url{https://aimath.org/WWN/hilbertschemes/hilbertschemes.pdf}. This URL is supplied provenance; no PDF access is claimed.
\item Supplied research summary, ``Generic points via elementary support factorization,'' attempt 1, refers to known elementary blocks with Hilbert functions \((1,4,3),(1,5,3),(1,6,3)\) but supplies no authors, titles, ideals, or complete theorem citations. Original remarks and literature fields are empty.
\end{itemize}
\end{document}
